\section*{Limitations}
\paragraph{Domain and population coverage.}
Each domain is represented by a particular benchmark and prompt format. ACPBench questions, for example, are not equivalent to unrestricted plan construction. The atlas uses one generation seed per prompt/model, and the annotated population is predominantly completed traces. Missing annotations and unfinished attempts limit generalization. Shared model families and task collections also mean that the seven checkpoints are not independent samples from a population of all reasoning models. We do not infer a scaling law or attribute differences causally to training lineage. Multiple benchmarks per domain, instruction paraphrases, and repeated generations would help separate task demand, format, and sampling variation.

\paragraph{Measurement validity and human comparisons.}
The labels are generated by Gemma4-31B-it. Comparing contextual and isolated annotations tests context sensitivity, not semantic correctness. The illustrative examples are selected without outcomes but are not a blinded human audit. A behavior can have different functions across tasks, and labels may partly track discourse style. Relative-position trajectories can obscure absolute duration and proximity to termination. Fixed-prefix analyses including unfinished attempts, event-aligned trajectories, and blinded coding of productive revision versus repetition would strengthen measurement. Marginal curves do not establish a universal sequence of reasoning states or faithfulness to internal computation. Cognitive-science connections are interpretive: we collect no human behavioral or neural data, and the two annotation families do not identify separate cognitive systems or brain networks. A matched human--model task-demand study could test functional similarities without assuming neural correspondence.

\paragraph{Intervention selection and inference.}
The main intervention experiments use Qwen3.5-9B. Archetypes were chosen through prior analyses and pilots, so A2 is not a prospectively selected optimum. Transfer cells share baselines, and the displayed paired intervals are exploratory rather than a new multiplicity-adjusted family. Equal vector norm does not match induced behavioral dose. Aggregate behavioral cosine lacks a joint uncertainty estimate, and shared baselines and response duration can contribute to similarity. Joint prompt resampling, duration-matched checks, and independently calibrated doses across models would strengthen comparison. Changes in behavior and performance do not identify mediation by a particular operation.

\paragraph{Dynamic mechanisms.}
The present analyses do not identify the dynamics underlying steering failures. One hypothesis is that some interventions induce repetitive dynamics even when activations remain statistically familiar. A subsequent study could test recurrence, responses to brief perturbations, and recovery of productive output, controlling intervention dose and decoding. Static activation density or projection onto a domain axis would not by itself establish an attractor or its causal role.

\paragraph{Reward learning and evaluation.}
Each RL initialization has one training seed. The earlier grading implementation supplied training rewards and informed archetype selection; updated evaluation grades cannot retroactively correct that objective. Replication with corrected rewards and multiple training seeds is needed before a training-general claim. Forced answering grants an answer cue and extra decoding, rather than reading latent knowledge directly. Comparisons should account for all answer-elicitation compute. Open-task preferences are judge-dependent and conditional on paired answer availability; neither judge agreement nor a zero-filled composite alone resolves that selection issue. Future evaluations should report completion and conditional quality jointly at matched budgets.

\section*{Ethical Considerations}
The atlas includes harmful-request and hazardous-knowledge benchmarks. Higher WMDP-Cyber accuracy measures greater hazardous knowledge, not safer behavior, and the present interventions are not validated as safety improvements. We report aggregate patterns without reproducing harmful instructions. Moral and creative quality judgments are limited operationalizations and should not be treated as objective measures of ethical worth or scientific value. Steering methods that alter reasoning or completion may also affect safety in ways not captured by the evaluated objectives. Any release should preserve source-dataset access conditions and licenses and avoid exposing identifying or harmful trace content unnecessarily.
